% =====================================================================
% fig_graph_construction.tex -- how a verification plan is constructed:
% name the things once, decide the measurements in small groups against
% those fixed names, assemble, then check the whole plan before anything
% runs.
%
% \input-able into main.tex, which loads adobe-techreport.sty (hence the
% adobered / adobeink / adobegray colours and the Source font family).
%
% REQUIRED IN THE MAIN PREAMBLE:
%     \usepackage{tikz}
%     \usetikzlibrary{arrows.meta,positioning,calc}
% Both lines are required, and the libraries MUST be loaded in the preamble --
% see the guard below for why a figure file cannot load them for itself.
% Every style is declared inside the tikzpicture's own option list, so all
% definitions are local to this picture: \input-ing this file twice, or
% alongside its companion figures, cannot clash.
%
% CONVENTIONS shared with fig_critic_pipeline.tex and fig_plan_example.tex:
%   adobegray!12 fill      = a learned component
%   white + adobeink line  = a deterministic component
%   adobered line          = the one check that precedes execution (used here
%                            for the whole-plan check and nothing else)
%   dashed gray            = something handed to the system, not produced by it
%
% CONTENT follows the subsection "Constructing the plan": the naming pass, the
% grouping of claims, per-group measurement writing against the fixed names,
% mechanical renaming, assembly, the whole-plan check, the re-ask of a failing
% group (format only), and the charged per-claim fallback. The figure is
% descriptive of the method only: no accuracy, no coverage, no comparison
% between configurations, and no quantity of any kind appears in it.
%
% Full-width boxes are 13.1cm, so the whole picture stays inside the report's
% text width unscaled -- do not add scale= or \resizebox. Compile, crop, render
% and width-check all three critic figures with
%     bash /Users/wenzhuox/diffphy_psc/techreport/figs/build_critic_figs.sh
% =====================================================================
% The TikZ libraries must be loaded in the PREAMBLE, not here: a figure
% environment forms a group, and TikZ defines a library's macros locally, so a
% \usetikzlibrary issued inside a figure does not survive to the next one. This
% guard turns a missing preamble line into a clear message instead of TikZ's
% confusing "You need to say \usetikzlibrary{calc}" further down.
\makeatletter
\@ifundefined{tikz@library@calc@loaded}{%
  \PackageError{fig_graph_construction}{%
    The preamble must contain\MessageBreak
    \string\usetikzlibrary{arrows.meta,positioning,calc}%
  }{Add that line after \string\usepackage{tikz}.}}{}
\makeatother

\begingroup
\providecommand{\gcsub}[1]{{\scriptsize\color{adobegray}#1}}

\begin{tikzpicture}[
    x=1cm, y=1cm,
    % --- shared box shape -------------------------------------------------
    gcbox/.style   = {rounded corners=2pt, line width=0.7pt, align=center,
                      inner sep=3.6pt, font=\sffamily\footnotesize,
                      text=adobeink},
    % --- LEARNED component: light gray fill --------------------------------
    gclearn/.style = {gcbox, draw=adobeink, fill=adobegray!12},
    % --- DETERMINISTIC component: white with an ink outline ----------------
    gcdet/.style   = {gcbox, draw=adobeink, fill=white},
    % --- accent: the check the assembled plan must pass before execution ---
    gcgate/.style  = {gcbox, draw=adobered, fill=white},
    % --- what is handed to the system --------------------------------------
    gcdata/.style  = {gcbox, draw=adobegray, fill=white, inner sep=3pt,
                      dash pattern=on 2pt off 1.6pt},
    gcline/.style  = {draw=adobeink, line width=0.7pt},
    gcarr/.style   = {gcline, -{Straight Barb[length=3.6pt,width=3.8pt]}},
    gcnote/.style  = {font=\sffamily\scriptsize, text=adobegray,
                      inner sep=0pt, align=left},
    gctag/.style   = {font=\sffamily\scriptsize, text=adobegray, fill=white,
                      inner xsep=3pt, inner ysep=1pt, align=center},
    gcsw/.style    = {rounded corners=1.2pt, line width=0.7pt, inner sep=0pt,
                      minimum width=11pt, minimum height=7.5pt},
  ]

  % Full-width rows are centred on x=0 and 13.1cm wide, so their left edge plus
  % inner sep sits at about x=-6.73; the return lane sits just outside that, and
  % the "format only" tag to its left is what sets the picture's left extent.

  % ---------------- the naming pass (learned) ------------------------------
  \node[gclearn, text width=13.1cm, anchor=north] (name) at (0,0)
    {\textbf{Name the things once}\\[1.5pt]
     \gcsub{a single pass over the whole prompt lists every distinct object or
     actor it mentions, and every\\ separate occurrence of an event, and fixes a
     name for each; nothing is measured in this pass}};

  \node[gcdata, text width=4.6cm, anchor=south]
    (prompt) at ($(name.north)+(0,0.40)$)
    {\textbf{Prompt and its claims}};

  % ---------------- grouping (deterministic) -------------------------------
  \node[gcdet, text width=13.1cm, anchor=north]
    (split) at ($(name.south)+(0,-0.40)$)
    {\textbf{Claims taken in small groups}\\[1.5pt]
     \gcsub{a few claims at a time, in the order they appear in the prompt}};

  % ---------------- per-group measurement writing (learned) ----------------
  \node[gclearn, text width=13.1cm, anchor=north]
    (build) at ($(split.south)+(0,-0.30)$)
    {\textbf{For each group: what would settle these claims}\\[1.5pt]
     \gcsub{the group is shown the whole prompt and the whole list of names, but
     only its own claims; it may\\ refer to any name on the list, and may not
     invent a new one}};

  % ---------------- mechanical renaming (deterministic) --------------------
  \node[gcdet, text width=13.1cm, anchor=north]
    (norm) at ($(build.south)+(0,-0.30)$)
    {\textbf{Names made group-local}\\[1.5pt]
     \gcsub{each group's measurements and checks are renumbered under that
     group's own label and its internal\\ references rewritten to match, so two
     groups cannot collide; claims keep the wording they were given}};

  % ---------------- assembly (deterministic) ------------------------------
  \node[gcdet, text width=13.1cm, anchor=north]
    (assemble) at ($(norm.south)+(0,-0.30)$)
    {\textbf{Groups assembled into one plan}};

  % ---------------- the check before execution (accented) ------------------
  \node[gcgate, text width=13.1cm, anchor=north]
    (gate) at ($(assemble.south)+(0,-0.30)$)
    {{\color{adobered}\textbf{The whole plan is checked, not each group}}\\[1.5pt]
     \gcsub{types, argument shapes and every cross-reference are checked over the
     assembled plan;\\ each complaint is traced back to the group it came from}};

  % ---------------- the two recoveries ------------------------------------
  \coordinate (fan) at ($(gate.south)+(0,-0.30)$);

  \node[gcdet, text width=5.8cm, anchor=north west]
    (retry) at ($(gate.south west)+(0,-0.76)$)
    {\textbf{Re-ask that group}\\[1.5pt]
     \gcsub{shown its own complaints and the rules\\ of the format --- never what
     to conclude}};

  \node[gcdet, text width=5.8cm, anchor=north east]
    (fall) at ($(gate.south east)+(0,-0.76)$)
    {\textbf{Otherwise, fall back claim by claim}\\[1.5pt]
     \gcsub{a group that still fails hands its claims\\ one at a time to the
     general verifier;\\ each is charged, none is dropped}};

  % ---------------- execution --------------------------------------------
  \coordinate (mrg) at ($(fall.south)+(0,-0.40)$);
  \node[gcdet, text width=13.1cm, anchor=north]
    (exec) at ($(0,0 |- mrg)+(0,-0.22)$)
    {\textbf{Execution}\\[1.5pt]
     \gcsub{a measurement requested more than once is executed once, and fixed
     rules turn the evidence\\ into the per-claim outcomes and the clip answer}};

  % ---------------- edges ------------------------------------------------
  \draw[gcarr] (prompt.south)   -- (name.north);
  \draw[gcarr] (name.south)     -- (split.north);
  \draw[gcarr] (split.south)    -- (build.north);
  \draw[gcarr] (build.south)    -- (norm.north);
  \draw[gcarr] (norm.south)     -- (assemble.north);
  \draw[gcarr] (assemble.south) -- (gate.north);

  % three outcomes of the check: the plan stands, re-ask a group, or fall back
  \draw[gcline] (gate.south) -- (fan);
  \draw[gcline] (retry.north |- fan) -- (fall.north |- fan);
  \draw[gcarr]  (retry.north |- fan) -- (retry.north);
  \draw[gcarr]  (fall.north  |- fan) -- (fall.north);
  \draw[gcarr]  (gate.south  |- fan) -- (exec.north);

  % the re-ask returns to the per-group pass, up the left-hand lane
  \draw[gcarr] (retry.west) -- ($(retry.west)+(-0.47,0)$) |- (build.west);

  \draw[gcarr] (fall.south) -- (fall.south |- exec.north);

  % labels drawn last, so the lines they sit on do not strike through them
  \node[gctag] at ($(gate.south |- fan)!0.85!(exec.north -| gate.south)$)
    {the plan stands};
  \node[gctag, anchor=south] at ($(retry.north)+(0,0.06)$) {a group failed};
  \node[gctag, anchor=south] at ($(fall.north)+(0,0.06)$) {still failing};
  \node[gctag, anchor=east] at ($(build.west)+(-0.47,0.40)$) {format only};

  % ---------------- inline legend ----------------------------------------
  \node[gcsw, draw=adobeink, fill=adobegray!12, anchor=north west]
    (swA) at ($(exec.south west)+(0,-0.44)$) {};
  \node[gcnote, right=3pt of swA] (lbA) {learned};

  \node[gcsw, draw=adobeink, fill=white, right=9pt of lbA] (swB) {};
  \node[gcnote, right=3pt of swB] (lbB) {deterministic};

  \node[gcsw, draw=adobered, fill=white, right=9pt of lbB] (swC) {};
  \node[gcnote, right=3pt of swC] (lbC) {checked before execution};

  \node[gcsw, draw=adobegray, fill=white, dash pattern=on 2pt off 1.6pt,
        right=9pt of lbC] (swD) {};
  \node[gcnote, right=3pt of swD] (lbD) {input};

\end{tikzpicture}
\endgroup
